\chapter*{Abstract}

The purpose of this thesis is to analyze the mathematical formulation behind Regev's quantum algorithm, explain its significance, practical functionality, and limitations, and present an example of its usefulness in a real-world scenario.

Specifically, we will begin by briefly explaining the idea behind Shor's algorithm, describing the classical and quantum subroutines, and analyzing its complexity. Following this, we will discuss certain mathematical concepts that provide the foundation for the generalizations and improvements introduced by Regev, such as the connection between finding periods in 1 dimension and determining bases for N-dimensional lattices,the probability of finding quadratic non-residues modulo N among primes, the Quantum Fourier Transform and its relation to the dual lattice and the Lenstra–Lenstra–Lovász algorithm.

Next, we will comprehensively explain Regev's algorithm and its theoretical complexity with respect to both gate and qubit requirements, while also mentioning certain improvements made since the publication of the original paper. Finally, we will demonstrate a practical implementation of the algorithm in Qiskit, illustrating an attempt to break a 2048-bit RSA key.